\chapter{The Spectral Theorem}

\begin{center}\itshape
A symmetric matrix is a list of numbers in disguise: choose the right axes and it only stretches. The same is true of the compact
self-adjoint operators, and the axes they choose are the vibrations of strings, the modes of noise and the harmonics of a drum. The Hermitian case was proved for matrices in Chapter~\ref{linear:ch:spectral}.
\end{center}

In 1910, the Dutch physicist Hendrik Lorentz came to G\"ottingen and lectured on an old problem of radiation. Inside a cavity, the
electromagnetic field vibrates at certain frequencies, the eigenvalues of the Laplacian with boundary conditions. Physicists needed to
know how many of them lie below a given frequency, and they believed, without proof, that the answer depends only on the volume of the
cavity and not on its shape. Lorentz asked the mathematicians in the audience to prove it. The story goes that Hilbert predicted he
would not live to see a proof. Less than two years later, his student Hermann Weyl proved it, with an idea that this chapter makes
precise in one dimension: the eigenvalues can be compared, one by one, through a variational principle, and squeezing a domain between
simpler ones squeezes its eigenvalues between theirs.

The setting had been built by Hilbert himself. Between 1904 and 1910 he wrote six papers on integral equations, in which he took
Fredholm's work, recast it as linear algebra in infinitely many variables, and proved that a symmetric kernel can be diagonalised. He
called the set of eigenvalues a \emph{spectrum}, by an analogy with the spectral lines of light that nobody then took literally.
Twenty years later the new quantum mechanics showed that the spectral lines of hydrogen \emph{are} the eigenvalues of an operator, and
Hilbert is reported to have said that he had named his theory after spectra without any presentiment that it would one day apply to the
actual spectra of physics.

This chapter proves the spectral theorem for compact self-adjoint operators, the variational characterisation of eigenvalues by
Courant and Fischer, and its two great applications: the vibrating string (Sturm--Liouville theory, with Weyl's law on an interval),
and Mercer's theorem on positive kernels, which expands the covariance of Brownian motion in its natural modes. It ends with a first
glimpse of bounded self-adjoint operators: their continuous functions, and their square roots. Throughout, $H$ is a complex Hilbert
space.

\section{Self-adjoint operators}

An operator $T\in\mathcal L(H)$ is self-adjoint when $\langle Tx,y\rangle=\langle x,Ty\rangle$ for all $x,y$. Then $\langle Tx,x\rangle=\langle x,Tx\rangle=\overline{\langle Tx,x\rangle}$ is real. The function $x\mapsto\langle Tx,x\rangle$ on the
unit sphere, the \emph{Rayleigh quotient}, carries all the information.

\begin{lemma}\label{lem:sanorm}
If $T$ is self-adjoint, then $\|T\|=\sup_{\|x\|=1}|\langle Tx,x\rangle|$.
\end{lemma}
\begin{proof}
Let $M$ be the supremum; clearly $M\le\|T\|$. Expanding, and using $\langle Ty,x\rangle=\overline{\langle Tx,y\rangle}$,
\[\langle T(x+y),x+y\rangle-\langle T(x-y),x-y\rangle=4\operatorname{Re}\langle Tx,y\rangle,\]
so $4\operatorname{Re}\langle Tx,y\rangle\le M(\|x+y\|^2+\|x-y\|^2)=2M(\|x\|^2+\|y\|^2)$ by the parallelogram law. For $\|x\|=1$ with $Tx\ne0$, take $y=\frac{Tx}{\|Tx\|}$: this gives $4\|Tx\|\le4M$.
\end{proof}

\begin{proposition}\label{prop:saspec}
Let $T$ be self-adjoint, $m=\inf_{\|x\|=1}\langle Tx,x\rangle$ and $M=\sup_{\|x\|=1}\langle Tx,x\rangle$. Then:
\begin{enumerate}
\item[(a)] the eigenvalues of $T$ are real, and eigenvectors for distinct eigenvalues are orthogonal;
\item[(b)] $\sigma(T)\subset[m,M]$, and $m,M\in\sigma(T)$;
\item[(c)] $\|T\|=\max(|m|,|M|)=r(T)$.
\end{enumerate}
\end{proposition}
\begin{proof}
(a) If $Tx=\lambda x$ with $\|x\|=1$, then $\lambda=\langle Tx,x\rangle\in\R$. If also $Ty=\mu y$ with $\mu\ne\lambda$, then $\lambda\langle x,y\rangle=\langle Tx,y\rangle=\langle x,Ty\rangle=\mu\langle x,y\rangle$, so $\langle x,y\rangle=0$.

(b) First, an operator $A$ with $\|Ax\|\ge c\|x\|$ for all $x$ ($c>0$) and $\ker A^*=\{0\}$ is invertible: it is injective, its range is closed (if $Ax_n\to y$, then $(x_n)$ is Cauchy, $x_n\to x$ and $Ax=y$), and the range is dense since its
orthogonal is $\ker A^*$ (Proposition~3.1); the inverse has norm at most $\frac1c$. Now let $\lambda=a+ib$ with $b\ne0$. For every $x$,
\[\|(T-\lambda)x\|\|x\|\ge|\operatorname{Im}\langle(T-\lambda)x,x\rangle|=|b|\|x\|^2,\]
and the same holds for $(T-\lambda)^*=T-\bar\lambda$. So $T-\lambda$ is invertible. If $\lambda>M$ is real, $\langle(\lambda-T)x,x\rangle\ge(\lambda-M)\|x\|^2$, so $\|(\lambda-T)x\|\ge(\lambda-M)\|x\|$, and $\lambda-T$ is self-adjoint:
it is invertible. The case $\lambda<m$ is symmetric. So $\sigma(T)\subset[m,M]$.

To see that $M\in\sigma(T)$, note that $B(x,y)=\langle(M-T)x,y\rangle$ is a positive Hermitian form, so it satisfies the Cauchy--Schwarz inequality $|B(x,y)|^2\le B(x,x)B(y,y)$ (Volume~0's proof applies word for word). With $y=(M-T)x$,
\[\|(M-T)x\|^4\le\langle(M-T)x,x\rangle\,\|M-T\|\,\|(M-T)x\|^2.\]
Choose unit vectors $x_n$ with $\langle Tx_n,x_n\rangle\to M$: then $(M-T)x_n\to0$. If $M-T$ had a bounded inverse, $x_n=(M-T)^{-1}(M-T)x_n$ would tend to $0$, which is absurd. So $M\in\sigma(T)$; applying this to $-T$ gives $m\in\sigma(T)$.

(c) By Lemma~\ref{lem:sanorm}, $\|T\|=\max(|m|,|M|)$, and by (b) this is the largest modulus of a point of $\sigma(T)$.
\end{proof}

So for a self-adjoint operator, the spectral radius formula of Theorem~3.6 is an equality at every step: $\|T^{2^k}\|=\|T\|^{2^k}$. An operator is \emph{positive}, written $T\ge0$, if $\langle Tx,x\rangle\ge0$ for all $x$; on a complex Hilbert space this forces $T$
to be self-adjoint (polarisation recovers $\langle Tx,y\rangle$ from the real numbers $\langle Tz,z\rangle$, which equal $\langle z,Tz\rangle$). By (b), $T\ge0$ if and only if $T$ is self-adjoint with $\sigma(T)\subset[0,\infty[$.

\section{The spectral theorem for compact operators}

\begin{lemma}\label{lem:eigenvalue}
A compact self-adjoint operator $K$ has $\|K\|$ or $-\|K\|$ as an eigenvalue.
\end{lemma}
\begin{proof}
If $K=0$ there is nothing to prove. By Lemma~\ref{lem:sanorm} there are unit vectors $x_n$ with $\langle Kx_n,x_n\rangle\to\lambda$, where $\lambda\in\{\|K\|,-\|K\|\}$. Then
\[\|Kx_n-\lambda x_n\|^2=\|Kx_n\|^2-2\lambda\langle Kx_n,x_n\rangle+\lambda^2\le2\lambda^2-2\lambda\langle Kx_n,x_n\rangle\to0.\]
By compactness, a subsequence has $Kx_n\to y$; then $\lambda x_n\to y$, so $\|y\|=|\lambda|\ne0$, and $Ky=\lim\lambda Kx_n=\lambda y$.
\end{proof}

\begin{theorem}[spectral theorem, Hilbert 1906]\label{thm:spectralcompact}
Let $K$ be a compact self-adjoint operator on $H$. There is an orthonormal family $(e_n)$, finite or countable, and non-zero real numbers $\lambda_n$, with $|\lambda_1|\ge|\lambda_2|\ge\cdots$ and $\lambda_n\to0$ if the family is
infinite, such that
\[Kx=\sum_n\lambda_n\langle x,e_n\rangle e_n\qquad\text{for all }x\in H.\]
Moreover $H=\ker K\oplus\overline{\operatorname{span}}(e_n)$; the non-zero eigenvalues of $K$ are exactly the $\lambda_n$, each of finite multiplicity; and if $H$ is separable, $H$ has a Hilbert basis of eigenvectors of $K$.
\end{theorem}
\begin{proof}
By Lemma~\ref{lem:eigenvalue} pick a unit $e_1$ with $Ke_1=\lambda_1e_1$, $|\lambda_1|=\|K\|$ (stop if $K=0$). The space $H_2=\{e_1\}^\perp$ is invariant under $K$: if $x\perp e_1$, then $\langle Kx,e_1\rangle=\langle x,Ke_1\rangle=0$. The restriction
of $K$ to $H_2$ is compact and self-adjoint, so it has a unit eigenvector $e_2$ with eigenvalue $\lambda_2$, $|\lambda_2|=\|K|_{H_2}\|\le|\lambda_1|$. Continue with $H_3=\{e_1,e_2\}^\perp$, and so on, stopping if some restriction
$K|_{H_{n+1}}$ vanishes.

If the process does not stop, $\lambda_n\to0$: otherwise $|\lambda_n|\ge\delta>0$ for all $n$, and $\|Ke_n-Ke_m\|^2=\lambda_n^2+\lambda_m^2\ge2\delta^2$, so $(Ke_n)$ has no convergent subsequence. For $x\in H$, the vector
$x-\sum_{k\le n}\langle x,e_k\rangle e_k$ lies in $H_{n+1}$ and has norm at most $\|x\|$, so
\[\Bigl\|Kx-\sum_{k\le n}\lambda_k\langle x,e_k\rangle e_k\Bigr\|\le|\lambda_{n+1}|\,\|x\|\to0\]
(and the difference is $0$ if the process stopped at $n$). This is the expansion. Let $F=\overline{\operatorname{span}}(e_n)$. If $x\in\ker K$, then $\lambda_n\langle x,e_n\rangle=\langle x,Ke_n\rangle=\langle Kx,e_n\rangle=0$, so
$x\perp F$; conversely if $x\perp F$, the expansion gives $Kx=0$. So $\ker K=F^\perp$ and $H=\ker K\oplus F$ (Corollary~2.2).

If $Kx=\mu x$ with $\mu\ne0$ and $x\ne0$, then $\mu\langle x,e_n\rangle=\langle Kx,e_n\rangle=\lambda_n\langle x,e_n\rangle$, so $x$ is orthogonal to the $e_n$ with $\lambda_n\ne\mu$, and $x\in F$ since $x=\mu^{-1}Kx$. Hence $\mu$
is one of the $\lambda_n$ and $x$ lies in the span of the $e_n$ with $\lambda_n=\mu$, a finite set since $\lambda_n\to0$. Finally, if $H$ is separable, so is $\ker K$, which has a countable Hilbert basis (Corollary~2.6); together
with the $e_n$ it is a Hilbert basis of $H$ made of eigenvectors.
\end{proof}

\begin{beauty}
The proof is an infinite version of the high-school recipe for a symmetric matrix: find the largest stretch, which is an eigenvalue,
remove its axis, and repeat. In finite dimension the maximum of the Rayleigh quotient on the sphere exists because the sphere is compact.
In infinite dimension the sphere is not compact (Theorem~1.3), and compactness has moved from the sphere to the operator: that is
exactly what is needed for the maximising sequence to converge.
\end{beauty}

\begin{think}
Does every compact operator have an eigenvalue? The Volterra operator of Example~3.7 is compact (Example~3.3: its kernel is $\mathbf 1_{y<x}$) and has none. What fails in Lemma~\ref{lem:eigenvalue}
without self-adjointness?
\end{think}

\begin{corollary}\label{cor:normal}
A compact normal operator $K$ on a separable Hilbert space has a Hilbert basis of eigenvectors, with complex eigenvalues tending to $0$.
\end{corollary}
\begin{proof}
$A=\frac12(K+K^*)$ and $B=\frac1{2i}(K-K^*)$ are compact and self-adjoint, $K=A+iB$, and $AB=BA$ because $K$ is normal. Each eigenspace $E$ of $A$ is invariant under $B$ (if $Ax=\lambda x$ then $ABx=BAx=\lambda Bx$), and $B|_E$ is compact
and self-adjoint, so by Theorem~\ref{thm:spectralcompact} applied in $E$ (finite-dimensional when $\lambda\ne0$, possibly not for $\lambda=0$), $E$ has a Hilbert basis of eigenvectors of $B$. These are eigenvectors of $A$ and $B$, hence of
$K$, and by Theorem~\ref{thm:spectralcompact} for $A$, $H$ is the orthogonal sum of the eigenspaces of $A$. Their union is the required basis. The eigenvalues of $K$ tend to $0$ by the argument in the proof of
Theorem~\ref{thm:spectralcompact}.
\end{proof}

\section{Min--max}

The largest eigenvalue of a compact positive operator maximises the Rayleigh quotient. The others are characterised in the same way, without
reference to the eigenvectors that come before them, and this makes them comparable from one operator to another.

\begin{theorem}[Courant--Fischer]\label{thm:minmax}
Let $K\ge0$ be compact on $H$, and let $\lambda_1\ge\lambda_2\ge\cdots\ge0$ be its non-zero eigenvalues repeated according to multiplicity, followed by zeros if there are finitely many. If $\dim H\ge n$, then
\begin{align*}
\lambda_n&=\max_{\dim V=n}\ \min_{x\in V,\,\|x\|=1}\langle Kx,x\rangle\\
&=\min_{\dim W\le n-1}\ \sup_{x\perp W,\,\|x\|=1}\langle Kx,x\rangle,
\end{align*}
where $V$ and $W$ run over subspaces of $H$.
\end{theorem}
\begin{proof}
Let $(e_k)$ be the orthonormal eigenvectors of Theorem~\ref{thm:spectralcompact}, ordered so that $Ke_k=\lambda_ke_k$ for the non-zero $\lambda_k$. For every $x$, $\langle Kx,x\rangle=\sum_k\lambda_k|\langle x,e_k\rangle|^2$.

If $x\perp e_1,\dots,e_{n-1}$, then $\langle Kx,x\rangle=\sum_{k\ge n}\lambda_k|\langle x,e_k\rangle|^2\le\lambda_n\|x\|^2$ by Bessel. Every subspace $V$ of dimension $n$ contains a unit vector orthogonal to $e_1,\dots,e_{n-1}$ (the
linear map $x\mapsto(\langle x,e_k\rangle)_{k<n}$ from $V$ to $\C^{n-1}$ has a non-zero kernel), so the minimum over $V$ is at most $\lambda_n$. Conversely, take $V$ spanned by those $e_k$, $k\le n$, which exist, completed if needed
by vectors of $\ker K$ (there are enough, since $\dim H\ge n$): for $x\in V$, $\langle Kx,x\rangle\ge\lambda_n\|x\|^2$. This proves the first formula, and the minimum over this $V$ is attained by compactness of its unit sphere.

For the second, $W=\operatorname{span}(e_k:k<n)$ gives a supremum at most $\lambda_n$ by the computation above. For any $W$ of dimension at most $n-1$, the $V$ just built meets $W^\perp$ in a non-zero vector, at which the
Rayleigh quotient is at least $\lambda_n$.
\end{proof}

\begin{corollary}\label{cor:weylpert}
If $K,L\ge0$ are compact, then $|\lambda_n(K)-\lambda_n(L)|\le\|K-L\|$ for every $n$. If $K\le L$, that is if $L-K\ge0$, then $\lambda_n(K)\le\lambda_n(L)$.
\end{corollary}
\begin{proof}
For unit $x$, $\langle Kx,x\rangle\le\langle Lx,x\rangle+\|K-L\|$. Take the minimum over $x\in V$, then the maximum over $V$, and use Theorem~\ref{thm:minmax}; exchange $K$ and $L$. If $K\le L$ the same argument works
without the error term.
\end{proof}

So the eigenvalues of positive compact operators move continuously, and monotonically, with the operator. For matrices this is
Weyl's inequality (1912); Exercise~\ref{ex:weylineq} gives the general form. Numerically, it says that computing eigenvalues of a
symmetric matrix is a well-conditioned problem, in sharp contrast to the Jordan block of Figure~3.1.

\section{Sturm--Liouville problems}

A string of length $1$, fixed at both ends, with a stiffness that varies along it, vibrates in modes $u(x)\cos(\omega t)$ where
\begin{equation}\label{eq:sl}
-u''+qu=\lambda u,\qquad u(0)=u(1)=0,
\end{equation}
with $\lambda=\omega^2$. Charles Sturm and Joseph Liouville studied such problems in 1836--1837, in a series of papers that Liouville
published in the journal he had just founded. The differential operator $u\mapsto-u''+qu$ is unbounded, but its inverse is a compact
integral operator, and the spectral theorem applies to the inverse. We assume $q\in C([0,1])$ real and $q\ge0$ (adding a constant to $q$
shifts every $\lambda$ by the same constant, so $q$ bounded below would do).

\begin{lemma}\label{lem:ivp}
For $x_0\in[0,1]$ and $a,b\in\C$ there is a unique $u\in C^2([0,1])$ with $u''=qu$, $u(x_0)=a$, $u'(x_0)=b$.
\end{lemma}
\begin{proof}
By Taylor's formula with integral remainder, these conditions are equivalent to $u=a+b(x-x_0)+Wu$, where $Wu(x)=\int_{x_0}^x(x-t)q(t)u(t)\dd t$ is a bounded operator on $C([0,1])$. By induction,
$|W^nu(x)|\le\frac{\|q\|_\infty^n|x-x_0|^{2n}}{(2n)!}\|u\|_\infty$, using $\int_0^s(s-t)t^{2n-2}\dd t=\frac{s^{2n}}{(2n-1)2n}$. So $\|W^n\|\le\frac{\|q\|_\infty^n}{(2n)!}$, the Neumann series $\sum W^n$ converges
(as in Lemma~3.5), $I-W$ is invertible, and the equation has exactly one continuous solution. It is $C^2$, and differentiating twice gives $u''=qu$.
\end{proof}

This is the Volterra phenomenon of Example~3.7 once more: a ``triangular'' integral operator has spectral radius $0$.

\begin{lemma}[Green's function]\label{lem:green}
Let $\varphi,\psi$ solve $\varphi''=q\varphi$, $\varphi(0)=0$, $\varphi'(0)=1$ and $\psi''=q\psi$, $\psi(1)=0$, $\psi'(1)=-1$. Then $c=\psi\varphi'-\psi'\varphi$ is a non-zero constant, and with
\[G(x,y)=\tfrac1c\,\varphi(\min(x,y))\,\psi(\max(x,y)),\]
for every $f\in C([0,1])$ the function $u=\mathcal Gf$, $\mathcal Gf(x)=\int_0^1G(x,y)f(y)\dd y$, is the unique $C^2$ solution of $-u''+qu=f$, $u(0)=u(1)=0$.
\end{lemma}
\begin{proof}
$c'=\psi\varphi''-\psi''\varphi=0$, so $c$ is constant, and evaluating at $x=1$ gives $c=\varphi(1)$. If $c=0$, then $\varphi(0)=\varphi(1)=0$ and $\varphi''=q\varphi$. Then
$0=\int_0^1\bar\varphi(\varphi''-q\varphi)=-\int_0^1(|\varphi'|^2+q|\varphi|^2)$ after an integration by parts, so $\varphi'=0$, contradicting $\varphi'(0)=1$. Now
$cu(x)=\psi(x)\int_0^x\varphi f+\varphi(x)\int_x^1\psi f$, so $u(0)=u(1)=0$, and differentiating,
\[cu'=\psi'\int_0^x\varphi f+\varphi'\int_x^1\psi f,\qquad cu''=q\,cu+(\psi'\varphi-\varphi'\psi)f=c\,(qu-f),\]
since the terms $\psi\varphi f-\varphi\psi f$ produced by differentiating the integrals cancel in $cu'$. So $-u''+qu=f$. If $w$ is the difference of two solutions, then $-w''+qw=0$ and $w(0)=w(1)=0$, and the integration by parts
above gives $w'=0$, so $w=0$.
\end{proof}

For $q=0$, $\varphi(x)=x$, $\psi(x)=1-x$, $c=1$, and $G(x,y)=\min(x,y)-xy$: a useful kernel for numerical methods, and the covariance of the Brownian bridge (Volume~3, Chapter~11). One more ingredient
is needed: enough test functions.

\begin{lemma}\label{lem:sines}
The functions $s_n(x)=\sqrt2\sin(n\pi x)$, $n\ge1$, form a Hilbert basis of $L^2([0,1])$.
\end{lemma}
\begin{proof}
They are orthonormal, since
\[2\sin(n\pi x)\sin(m\pi x)=\cos((n-m)\pi x)-\cos((n+m)\pi x)\]
integrates to $\delta_{nm}$ over $[0,1]$. Let $f\in L^2([0,1])$ be orthogonal to all of them, and extend it to an odd function $\tilde f$ on $[-1,1]$. Then
$\int_{-1}^1\tilde f(x)\sin(n\pi x)\dd x=2\int_0^1f(x)\sin(n\pi x)\dd x=0$, and $\int_{-1}^1\tilde f(x)\cos(n\pi x)\dd x=0$ because the integrand is odd. So $\tilde f$ is orthogonal to every $e^{in\pi x}$, and $t\mapsto\tilde f(\frac t\pi)$ is
orthogonal to every $e^{int}$ in $L^2(\mathbb T)$; by Theorem~2.7 it vanishes, so $f=0$.
\end{proof}

\begin{theorem}[Sturm--Liouville]\label{thm:sturm}
Let $q\in C([0,1])$ be real with $q\ge0$. There is a Hilbert basis $(e_n)_{n\ge1}$ of $L^2([0,1])$ of real functions $e_n\in C^2([0,1])$ with $e_n(0)=e_n(1)=0$ and $-e_n''+qe_n=\lambda_ne_n$, where
\[0<\lambda_1<\lambda_2<\cdots,\qquad \pi^2n^2\le\lambda_n\le\pi^2n^2+\max q.\]
Every non-zero $C^2$ solution of \eqref{eq:sl} is a multiple of some $e_n$, with $\lambda=\lambda_n$.
\end{theorem}
\begin{proof}
\emph{The inverse operator.} $G$ is continuous, real and symmetric, so $\mathcal G$ is a compact self-adjoint operator on $L^2([0,1])$ (Example~3.3). It maps $L^2$ into $C([0,1])$, since
$|\mathcal Gf(x)-\mathcal Gf(x')|\le\|G(x,\cdot)-G(x',\cdot)\|_2\|f\|_2$ and $G$ is uniformly continuous. For $f$ continuous and $u=\mathcal Gf$, integration by parts gives
\[\langle\mathcal Gf,f\rangle=\int_0^1u\,\overline{(-u''+qu)}=\int_0^1\bigl(|u'|^2+q|u|^2\bigr)\ge0,\]
and by density of $C([0,1])$ in $L^2$ (Volume~6, Chapter~9) and continuity of $f\mapsto\langle\mathcal Gf,f\rangle$, $\mathcal G\ge0$. It is injective: if $\mathcal Gf=0$, then for every finite combination $h$ of the $s_n$ of
Lemma~\ref{lem:sines}, which is $C^2$ with $h(0)=h(1)=0$, uniqueness in Lemma~\ref{lem:green} gives $\mathcal G(-h''+qh)=h$, so $\langle f,h\rangle=\langle\mathcal Gf,-h''+qh\rangle=0$; these $h$ are dense, so $f=0$.

\emph{Eigenfunctions.} By Theorem~\ref{thm:spectralcompact}, with $\ker\mathcal G=\{0\}$, $L^2([0,1])$ has a Hilbert basis $(e_n)$ with $\mathcal Ge_n=\mu_ne_n$, $\mu_n>0$ (positivity and injectivity), $\mu_n\to0$. Each $e_n=\mathcal G(\mu_n^{-1}e_n)$ is continuous,
hence $C^2$ by Lemma~\ref{lem:green}, with $e_n(0)=e_n(1)=0$ and $-e_n''+qe_n=\lambda_ne_n$, $\lambda_n=\mu_n^{-1}$. Conversely, a non-zero $C^2$ solution $u$ of \eqref{eq:sl} satisfies $u=\mathcal G(\lambda u)$, so $\lambda\ne0$ and $u$ is an
eigenvector of $\mathcal G$ for $\lambda^{-1}$: $\lambda$ is some $\lambda_n$. Two solutions $u,v$ of \eqref{eq:sl} for the same $\lambda$ are proportional, since $v'(0)u-u'(0)v$ solves the same equation with zero initial data
at $0$ and vanishes by Lemma~\ref{lem:ivp}. So every eigenvalue is simple, the $\lambda_n$ are distinct and can be listed increasingly, and as real and imaginary parts of a solution are solutions, each $e_n$ can be taken real.

\emph{Bounds.} For $u\in C^2$ with $u(0)=u(1)=0$, let $Q(u)=\int_0^1(|u'|^2+q|u|^2)=\langle-u''+qu,u\rangle$. Two integrations by parts give $\langle-u''+qu,e_k\rangle=\langle u,-e_k''+qe_k\rangle=\lambda_k\langle u,e_k\rangle$, so by Parseval
\[Q(u)=\sum_k\lambda_k|\langle u,e_k\rangle|^2.\]
The counting argument in the proof of Theorem~\ref{thm:minmax}, with the inequalities reversed because these $\lambda_n$ increase rather than decrease, now gives
\[\lambda_n=\min_V\ \max_{u\in V,\,\|u\|_2=1}Q(u),\]
the minimum over $n$-dimensional spaces $V$ of $C^2$ functions vanishing at $0$ and $1$, attained at $V=\operatorname{span}(e_1,\dots,e_n)$. Write $Q_q$ to show the dependence on $q$. Since $Q_0(u)\le Q_q(u)\le Q_0(u)+\max q\,\|u\|_2^2$,
the eigenvalues for $q$ lie between those for $0$ and those for $0$ plus $\max q$. For $q=0$, the solutions of $u''=-\lambda u$ with $u(0)=0$ are the multiples of $\sin(\sqrt\lambda x)$ (as $\lambda>0$), and $u(1)=0$ means
$\sqrt\lambda\in\pi\N^*$: the eigenvalues are the $\pi^2n^2$.
\end{proof}

\begin{figure}[ht]
\centering
\fig{ch04-sturm}
\caption{The first four eigenfunctions of $-u''+qu=\lambda u$ on $[0,1]$ with $q(x)=60x$ (solid), against the sines of $q=0$ (dashed). Where $q$ is large the string is stiffer and the modes are
pushed to the left; the $n$-th mode still has $n-1$ zeros inside, as Sturm's oscillation theorem predicts. The eigenvalues are $\lambda_n\approx36.2,\ 70.4,\ 119.5,\ 188.3$, against $\pi^2n^2\approx9.9,\ 39.5,\ 88.8,\ 157.9$.}
\label{fig:sturm}
\end{figure}

\begin{corollary}[Weyl's law on an interval]\label{cor:weyllaw}
Let $N(\lambda)$ be the number of eigenvalues of \eqref{eq:sl} at most $\lambda$. Then $N(\lambda)\sim\frac{\sqrt\lambda}\pi$ as $\lambda\to\infty$. On an interval of length $\ell$, with $q=0$, the eigenvalues are $\frac{\pi^2n^2}{\ell^2}$
and $N(\lambda)\sim\frac{\ell\sqrt\lambda}\pi$.
\end{corollary}
\begin{proof}
By Theorem~\ref{thm:sturm}, $\lambda_n\le\lambda$ implies $n\le\frac{\sqrt\lambda}\pi$, and $\pi^2n^2+\max q\le\lambda$ implies $\lambda_n\le\lambda$. So $\frac{\sqrt{\lambda-\max q}}\pi-1\le N(\lambda)\le\frac{\sqrt\lambda}\pi$. For the
interval $[0,\ell]$, $u(x)=v(\frac x\ell)$ turns $-u''=\lambda u$ into $-v''=\ell^2\lambda v$.
\end{proof}

\begin{think}
Can you hear the length of a string? The corollary says yes: $\ell=\lim\pi\frac{N(\lambda)}{\sqrt\lambda}$, and the potential $q$ is inaudible at this order. Can you hear $q$ itself? Partly. Ambartsumian proved in 1929 that if the
eigenvalues with the Neumann conditions $u'(0)=u'(1)=0$ are exactly the $\pi^2n^2$, $n\ge0$, then $q=0$; Borg showed in 1946 that two spectra, for two different boundary conditions, determine $q$, while one spectrum in general does not.
\end{think}

\begin{remark}
In 1966 Mark Kac asked the same question for a drum: does the spectrum of the Laplacian on a plane domain determine its shape? Weyl's law says it determines the area, and finer asymptotics give the perimeter and the
number of holes. In 1992 Gordon, Webb and Wolpert answered no, with two different polygons with exactly the same eigenvalues.
\end{remark}

\section{Hilbert--Schmidt operators and Mercer's theorem}

\begin{proposition}\label{prop:hs}
Let $T\in\mathcal L(H)$, with $H$ separable. The quantity $\|T\|_{HS}^2=\sum_n\|Te_n\|^2\in[0,\infty]$ does not depend on the Hilbert basis $(e_n)$, and $\|T^*\|_{HS}=\|T\|_{HS}$. If it is finite, $T$ is called
\emph{Hilbert--Schmidt}; then $\|T\|\le\|T\|_{HS}$ and $T$ is compact. On $H=L^2(\mu)$, $\mu$ $\sigma$-finite with $L^2(\mu)$ separable, the Hilbert--Schmidt operators are exactly the integral operators of
Example~3.3 with kernel $k\in L^2(\mu\otimes\mu)$, and $\|T\|_{HS}=\|k\|_{L^2}$.
\end{proposition}
\begin{proof}
For two Hilbert bases $(e_n)$ and $(f_m)$, Parseval and Tonelli for series give
\[\sum_n\|Te_n\|^2=\sum_n\sum_m|\langle Te_n,f_m\rangle|^2=\sum_m\sum_n|\langle e_n,T^*f_m\rangle|^2=\sum_m\|T^*f_m\|^2.\]
The right side does not involve $(e_n)$, and the left side does not involve $(f_m)$, so both are independent of the bases and $\|T\|_{HS}=\|T^*\|_{HS}$. If $\|x\|=1$, complete $x$ to a Hilbert basis: $\|Tx\|^2\le\|T\|_{HS}^2$.
With $P_N$ the projection on $\operatorname{span}(e_n:n\le N)$, $\|T-TP_N\|^2\le\|T-TP_N\|_{HS}^2=\sum_{n>N}\|Te_n\|^2\to0$, so $T$ is a limit of finite-rank operators, hence compact (Proposition~3.2).

On $L^2(\mu)$, the functions $e_m\otimes\bar e_n:(x,y)\mapsto e_m(x)\overline{e_n(y)}$ are orthonormal in $L^2(\mu\otimes\mu)$ by Fubini. They form a Hilbert basis: write $T_k$ for the integral operator with kernel $k$; then
$\langle k,e_m\otimes\bar e_n\rangle=\langle T_ke_n,e_m\rangle$, so if $k$ is orthogonal to all of them, $T_k=0$, hence $\int\!\!\int k(x,y)\mathbf 1_B(y)\mathbf 1_A(x)=0$ for all sets $A,B$ of finite measure, and $k=0$ almost
everywhere, since a function whose integral vanishes on all rectangles of finite measure vanishes (Volume~6, Chapter~8: the rectangles generate the product $\sigma$-algebra). Now $\|T_k\|_{HS}^2=\sum_{m,n}|\langle T_ke_n,e_m\rangle|^2=\sum_{m,n}|\langle k,e_m\otimes\bar e_n\rangle|^2=\|k\|_2^2$.
Conversely, if $T$ is Hilbert--Schmidt, $k=\sum_{m,n}\langle Te_n,e_m\rangle\,e_m\otimes\bar e_n$ converges in $L^2(\mu\otimes\mu)$, and $T_k$ and $T$ agree on every $e_n$, so $T=T_k$.
\end{proof}

The kernel of a Hilbert--Schmidt operator is only defined almost everywhere, so its diagonal $k(x,x)$ means nothing in general. For continuous positive kernels it does, and it carries the trace.

\begin{theorem}[Mercer, 1909]\label{thm:mercer}
Let $I=[a,b]$ and let $k\in C(I\times I)$ be a \emph{positive kernel}: $\overline{k(x,y)}=k(y,x)$, and $\sum_{i,j}\bar c_ic_jk(x_i,x_j)\ge0$ for all finite families of points $x_i\in I$ and $c_i\in\C$. Let $T$ be the integral
operator with kernel $k$ on $L^2(I)$, with non-zero eigenvalues $\lambda_n$ and orthonormal eigenvectors $e_n$ as in Theorem~\ref{thm:spectralcompact}. Then $\lambda_n>0$, each $e_n$ is continuous, and
\[k(x,y)=\sum_n\lambda_ne_n(x)\overline{e_n(y)},\]
the series converging absolutely and uniformly on $I\times I$. In particular $\sum_n\lambda_n=\int_Ik(x,x)\dd x$.
\end{theorem}
\begin{proof}
\emph{Step 1: $T\ge0$, and the $e_n$ are continuous.} For $f$ continuous, the integral
\[\langle Tf,f\rangle=\int_I\!\int_Ik(x,y)f(y)\overline{f(x)}\dd x\dd y\]
is a limit, over regular grids of step $h$, of the Riemann sums
\[h^2\sum_{i,j}k(x_i,x_j)f(x_j)\overline{f(x_i)}\ge0;\] by density, $T\ge0$, so $\lambda_n>0$. As in the proof of Theorem~\ref{thm:sturm}, $T$ maps $L^2$ into $C(I)$, so $e_n=\lambda_n^{-1}Te_n$ is continuous.

\emph{Step 2: a bound on the diagonal.} Let $k_N(x,y)=k(x,y)-\sum_{n\le N}\lambda_ne_n(x)\overline{e_n(y)}$, a continuous kernel. Its operator $T-\sum_{n\le N}\lambda_n\langle\cdot,e_n\rangle e_n$ is positive, by the expansion of
Theorem~\ref{thm:spectralcompact}. A continuous kernel $\kappa$ whose operator is positive has $\kappa(x,x)\ge0$: $\kappa(x_0,x_0)$ is real, and if it were $<0$, then $\operatorname{Re}\kappa<\frac12\kappa(x_0,x_0)$ on some
$U\times U$, $U$ an interval around $x_0$, and with $f=\mathbf 1_U$ the real number $\langle T_\kappa f,f\rangle=\int_U\int_U\kappa$ would be negative. So $k_N(x,x)\ge0$, that is,
\[\sum_{n\le N}\lambda_n|e_n(x)|^2\le k(x,x)\le C:=\max_Ik(x,x)\qquad\text{for all }x\text{ and }N.\]

\emph{Step 3: pointwise convergence.} By Cauchy--Schwarz, for $M<N$,
\begin{equation}\label{eq:mercercs}
\begin{aligned}\sum_{n=M}^N\lambda_n|e_n(x)e_n(y)|&\le\Bigl(\sum_{n=M}^N\lambda_n|e_n(x)|^2\Bigr)^{1/2}\eqbr\cdot\Bigl(\sum_{n=M}^N\lambda_n|e_n(y)|^2\Bigr)^{1/2}.\end{aligned}
\end{equation}
For fixed $x$ the first factor tends to $0$ as $M\to\infty$ (Step 2) and the second is at most $\sqrt C$, so $\tilde k(x,y)=\sum_n\lambda_ne_n(x)\overline{e_n(y)}$ converges uniformly in $y$, and is continuous in $y$. For
$f\in L^2(I)$, the same inequality shows that $\sum_n\lambda_n\langle f,e_n\rangle e_n(x)$ converges uniformly in $x$; its limit in $L^2$ is $Tf$ (Theorem~\ref{thm:spectralcompact}), and both $Tf$ and the uniform limit
are continuous, so they agree everywhere. Integrating term by term (uniform convergence in $y$),
\[\int_I\tilde k(x,y)f(y)\dd y=\sum_n\lambda_ne_n(x)\langle f,e_n\rangle=Tf(x)=\int_Ik(x,y)f(y)\dd y.\]
As this holds for all $f$, $\tilde k(x,\cdot)=k(x,\cdot)$ almost everywhere, hence everywhere by continuity. In particular $k(x,x)=\sum_n\lambda_n|e_n(x)|^2$.

\emph{Step 4: uniform convergence.} The partial sums $\sum_{n\le N}\lambda_n|e_n(x)|^2$ are continuous and increase to the continuous function $k(x,x)$ on the compact $I$, so the convergence is uniform by Dini's theorem (Volume~6,
Exercise~4.1, or directly: the closed sets $\{x:k(x,x)-\sum_{n\le N}\lambda_n|e_n(x)|^2\ge\varepsilon\}$ decrease and have empty intersection, so one of them is empty). Then the first factor in \eqref{eq:mercercs} tends to
$0$ uniformly in $x$, and the series converges absolutely and uniformly on $I\times I$. Integrating the diagonal, $\int_Ik(x,x)\dd x=\sum_n\lambda_n\int_I|e_n|^2=\sum_n\lambda_n$.
\end{proof}

\begin{example}[the covariance of Brownian motion]\label{ex:minkernel}
The kernel $k(s,t)=\min(s,t)$ on $[0,1]$ is positive, because $\min(s,t)=\int_0^1\mathbf 1_{[0,s]}\mathbf 1_{[0,t]}$ gives $\sum\bar c_ic_j\min(t_i,t_j)=\int_0^1\bigl|\sum_jc_j\mathbf 1_{[0,t_j]}\bigr|^2\ge0$. If $Tf=\lambda f$ with
$\lambda\ne0$, then $f$ is continuous and
\[\lambda f(s)=\int_0^st\,f(t)\dd t+s\int_s^1f(t)\dd t,\]
so $\lambda f(0)=0$, $\lambda f'(s)=\int_s^1f$, $\lambda f'(1)=0$ and $\lambda f''=-f$. With $\lambda>0$, $f$ is a multiple of $\sin(s/\sqrt\lambda)$, and $f'(1)=0$ forces $\frac1{\sqrt\lambda}=(n-\frac12)\pi$. So
\[\lambda_n=\frac1{(n-\frac12)^2\pi^2},\qquad e_n(s)=\sqrt2\sin\bigl((n-\tfrac12)\pi s\bigr),\qquad n\ge1,\]
and Mercer's theorem gives, uniformly on $[0,1]^2$,
\[\min(s,t)=\sum_{n\ge1}\frac{2\sin((n-\frac12)\pi s)\sin((n-\frac12)\pi t)}{(n-\frac12)^2\pi^2}.\]
At $s=t=1$: $1=\frac8{\pi^2}\sum_{n\ge1}\frac1{(2n-1)^2}$, that is $\sum_{n\ge1}\frac1{(2n-1)^2}=\frac{\pi^2}8$, the odd half of the Basel problem.
\end{example}

Here is why this example matters. Brownian motion $(B_t)_{t\in[0,1]}$ (Volume~3, Chapter~11) is a centred Gaussian process with $\operatorname{Cov}(B_s,B_t)=\min(s,t)$. The variables $\xi_n=\lambda_n^{-1/2}\int_0^1B_te_n(t)\dd t$ are
jointly Gaussian with
\[\operatorname{Cov}(\xi_n,\xi_m)=\frac{\int_0^1\!\int_0^1\min(s,t)e_n(s)e_m(t)\dd s\dd t}{\sqrt{\lambda_n\lambda_m}}=\frac{\langle Te_m,e_n\rangle}{\sqrt{\lambda_n\lambda_m}}=\delta_{nm},\]
so they are independent standard Gaussians, and Brownian motion should be the random series $\sum_n\xi_n\sqrt{\lambda_n}e_n(t)$, a sum of sine waves with independent random amplitudes. This is the
\emph{Karhunen--Lo\`eve expansion}, the principal component analysis of a random function: the first mode carries $\lambda_1/\sum\lambda_n=\frac{8}{\pi^2}\approx81\%$ of the total variance $\int_0^1\operatorname{Var}B_t\dd t=\frac12$.
Exercise~\ref{ex:kl} proves that the series has the right covariances, and Volume~11 proves that it converges uniformly, almost surely, to a continuous path.

\begin{figure}[ht]
\centering
\fig{ch04-kl}
\caption{Left: the first four Karhunen--Lo\`eve modes $\sqrt{\lambda_n}e_n(t)$ of Brownian motion. Right: one sample of the partial sums $\sum_{n\le N}\xi_n\sqrt{\lambda_n}e_n(t)$ for $N=3$, $20$ and $500$, with the same random
amplitudes $\xi_n$: the low modes give the overall shape, the high ones the roughness.}
\label{fig:kl}
\end{figure}

\section{The continuous functional calculus}

For a compact self-adjoint operator, the spectral theorem defines
\[
f(K)=\sum_nf(\lambda_n)\langle\cdot,e_n\rangle e_n+f(0)P_{\ker K}
\]
 for any bounded function $f$ on the spectrum. For a general bounded self-adjoint operator, such as
multiplication by $x$ on $L^2([0,1])$, there may be no eigenvectors at all. Yet continuous functions of it still make sense, by an argument that uses nothing but Proposition~\ref{prop:saspec} and Weierstrass.

\begin{theorem}[continuous functional calculus]\label{thm:funcalc}
Let $T$ be a bounded self-adjoint operator on $H$. There is a unique linear map $f\mapsto f(T)$ from $C(\sigma(T))$ to $\mathcal L(H)$ which is multiplicative, sends $1$ to $I$ and the identity function to $T$, and is
isometric: $\|f(T)\|=\max_{\sigma(T)}|f|$. Moreover $f(T)^*=\bar f(T)$, $f(T)$ commutes with every operator that commutes with $T$, and $f(T)\ge0$ when $f\ge0$.
\end{theorem}
\begin{proof}
\emph{Polynomials.} For a polynomial $p$, $\sigma(p(T))=p(\sigma(T))$. Indeed, for $\mu\in\C$ and $p$ non-constant, factor $p(z)-\mu=c\prod_i(z-z_i)$, so $p(T)-\mu=c\prod_i(T-z_i)$ with commuting factors. A product $AB=BA$ of
commuting operators is invertible if and only if both factors are ($A$ then has the right inverse $B(AB)^{-1}$ and the left inverse $(BA)^{-1}B$). So $\mu\in\sigma(p(T))$ iff some $z_i\in\sigma(T)$ iff $\mu\in p(\sigma(T))$.
Write $\bar p$ for the polynomial with conjugate coefficients; as $T^*=T$, $p(T)^*=\bar p(T)$, and $(\bar pp)(T)$ is self-adjoint. By Proposition~3.1 and Proposition~\ref{prop:saspec}(c),
\[\|p(T)\|^2=\|p(T)^*p(T)\|=r\bigl((\bar pp)(T)\bigr)=\max_{t\in\sigma(T)}|\bar p(t)p(t)|=\max_{\sigma(T)}|p|^2,\]
using $\bar p(t)=\overline{p(t)}$ for real $t$, as $\sigma(T)\subset\R$.

\emph{Extension.} So $p\mapsto p(T)$ is well defined and isometric on polynomials viewed as functions on the compact set $\sigma(T)\subset\R$ (two polynomials equal on $\sigma(T)$ give the same operator). Polynomials are dense in
$C(\sigma(T))$ by the Weierstrass theorem (Volume~6, Theorem~4.19), and $\mathcal L(H)$ is complete, so the map extends uniquely to an isometry on $C(\sigma(T))$; linearity, multiplicativity, $f(T)^*=\bar f(T)$ and commutation
pass to the limit. Uniqueness: any map with these properties agrees with this one on polynomials and is continuous.

\emph{Positivity.} If $f\ge0$, let $g=\sqrt f$, real and continuous: $f(T)=g(T)^2=g(T)^*g(T)$, and $\langle f(T)x,x\rangle=\|g(T)x\|^2\ge0$.
\end{proof}

\begin{corollary}[square roots]\label{cor:sqrt}
Every positive operator $A$ has a unique positive square root: an operator $B\ge0$ with $B^2=A$.
\end{corollary}
\begin{proof}
Since $\sigma(A)\subset[0,\infty[$ (Proposition~\ref{prop:saspec}), $C=\sqrt{A}$, the function $\sqrt t$ applied to $A$, is positive and $C^2=A$. Let $B\ge0$ with $B^2=A$. Then $B$ commutes with $A=B^2$, hence with $C$
(Theorem~\ref{thm:funcalc}). So $(B-C)(B+C)=B^2-C^2=0$, and
\[(B-C)B(B-C)+(B-C)C(B-C)=(B-C)(B+C)(B-C)=0.\]
Both terms are positive ($\langle(B-C)B(B-C)x,x\rangle=\langle By,y\rangle$ with $y=(B-C)x$), so both vanish, and their difference $(B-C)^3$ vanishes. For the self-adjoint $S=B-C$, Proposition~3.1 gives $\|S^2\|=\|S\|^2$, so
$\|S\|^4=\|S^4\|=\|S\cdot S^3\|=0$ and $B=C$.
\end{proof}

\begin{remark}
The full spectral theorem says that every bounded self-adjoint operator is a multiplication operator in disguise: there is a measure space $(X,\mu)$, a bounded real function $\varphi$ on $X$, and a unitary
$U:H\to L^2(\mu)$ with $UTU^{-1}f=\varphi f$. Equivalently, there is a projection-valued measure $E$ on $\sigma(T)$ with $T=\int\lambda\dd E(\lambda)$, which allows $f(T)$ for every bounded Borel $f$. The proof
combines Theorem~\ref{thm:funcalc} with the Riesz--Markov theorem (Theorem~1.13); von Neumann extended it to unbounded operators in 1929, to give quantum mechanics its mathematics. See Reed and Simon, or Rudin, below.
\end{remark}

\section{Exercises}

\begin{exercise}[$\star\star$]\label{ex:kl}
\emph{The Karhunen--Lo\`eve expansion of Brownian motion.} Let $(\xi_n)_{n\ge1}$ be independent $N(0,1)$ variables, and $\lambda_n,e_n$ as in Example~\ref{ex:minkernel}. (a) Show that for each $t\in[0,1]$ the series
$X_t=\sum_n\xi_n\sqrt{\lambda_n}e_n(t)$ converges in $L^2(P)$. (b) Show that $X_0=0$ and $\operatorname{Cov}(X_s,X_t)=\min(s,t)$, and that each $X_t$ is Gaussian. (c) Find the corresponding expansion of the Brownian bridge,
whose covariance is $\min(s,t)-st$ (Volume~3, Chapter~11).
\end{exercise}

\begin{exercise}[$\star$]\label{ex:toeplitz}
\emph{Convolution kernels.} Let $g\in L^2(\mathbb T)$, with $\mathbb T=\R/\Z$ here, and $Tf(x)=\int_0^1g(x-y)f(y)\dd y$ on $L^2(\mathbb T)$. (a) Show that $e_n(x)=e^{2\pi inx}$ is an eigenvector of $T$, with eigenvalue the Fourier coefficient
$\hat g(n)=\int_0^1g(x)e^{-2\pi inx}\dd x$. (b) When is $T$ self-adjoint? positive? (c) Apply Mercer's theorem to the Poisson kernel $P_r(x)=\sum_{n\in\Z}r^{|n|}e^{2\pi inx}$, $0<r<1$, and check the trace formula.
\end{exercise}

\begin{exercise}[$\star\star$]\label{ex:weylineq}
\emph{Weyl's inequalities.} Let $K,L\ge0$ be compact. Show that $\lambda_{m+n-1}(K+L)\le\lambda_m(K)+\lambda_n(L)$ for all $m,n\ge1$. Deduce that if $F\ge0$ has rank $r$, then $\lambda_n(K)\le\lambda_n(K+F)$ and $\lambda_{n+r}(K+F)\le\lambda_n(K)$: a
positive perturbation of rank $r$ raises the eigenvalues, but never moves one past the eigenvalue $r$ places above it.
\end{exercise}

\begin{exercise}[$\star$]\label{ex:graphlaplacian}
\emph{The Laplacian of a graph.} Let $G$ be a finite graph with vertex set $V$ and edge set $E$, and $L=D-A$ on $\C^V$, where $D$ is the diagonal matrix of degrees and $A$ the adjacency matrix. (a) Show that
$\langle Lx,x\rangle=\sum_{\{i,j\}\in E}|x_i-x_j|^2$, that $L\ge0$, and that $\ker L$ is the space of functions constant on each connected component. (b) If $G$ is connected, show that the second smallest eigenvalue of $L$ is
$\mu_2=\min\bigl\{\frac{\sum_E|x_i-x_j|^2}{\sum_V|x_i|^2}:x\ne0,\ \sum_Vx_i=0\bigr\}$. (c) Compute the eigenvalues of the cycle $C_n$, and show that $\mu_2\sim\frac{4\pi^2}{n^2}$.
\end{exercise}

\begin{exercise}[$\star\star$]\label{ex:polar}
\emph{Polar decomposition.} Let $T\in\mathcal L(H)$ and $|T|=\sqrt{T^*T}$ (Corollary~\ref{cor:sqrt}). (a) Show that $\||T|x\|=\|Tx\|$ for all $x$. (b) Show that there is a unique $U\in\mathcal L(H)$ with $T=U|T|$ and $U=0$ on
$\ker T$, and that $U$ is isometric on $(\ker T)^\perp$. (c) Show that $U$ is unitary if $T$ is invertible, and compute $|S|$ and $U$ for the shift $S$ of Chapter~16.
\end{exercise}

\begin{exercise}[$\star$]\label{ex:basel}
\emph{Basel from a string.} Let $G(x,y)=\min(x,y)-xy$ on $[0,1]^2$. (a) Show that $G$ is a positive kernel. (b) Using Theorem~\ref{thm:sturm} with $q=0$ and Mercer's theorem, prove $\sum_{n\ge1}\frac1{n^2}=\frac{\pi^2}6$.
(c) Using Proposition~\ref{prop:hs}, prove $\sum_{n\ge1}\frac1{n^4}=\frac{\pi^4}{90}$.
\end{exercise}

\begin{exercise}[$\star\star$]\label{ex:volterranorm}
\emph{The norm of the Volterra operator.} Let $Vf(x)=\int_0^xf$ on $L^2([0,1])$. Compute $V^*$ and $V^*V$, find the eigenvalues of $V^*V$ with Example~\ref{ex:minkernel}, and deduce that $\|V\|=\frac2\pi$. (The operator $V$ has no
eigenvalue, but $V^*V$ has plenty: these \emph{singular values} are the right tool for non-normal compact operators.)
\end{exercise}

\section{Hints}

\begin{hint}{ex:kl}
The terms are independent and centred, so they are orthogonal in $L^2(P)$. Mercer's theorem on the diagonal gives $\sum_n\lambda_ne_n(t)^2=t$. For (c), the kernel is the Green's function of $-u''$.
\end{hint}
\begin{hint}{ex:toeplitz}
Substitute $u=x-y$. Since the $e_n$ form a Hilbert basis (Theorem~2.7, rescaled), $T$ is diagonal in it, and $\langle Tf,f\rangle=\sum_n\hat g(n)|\langle f,e_n\rangle|^2$.
\end{hint}
\begin{hint}{ex:weylineq}
Use the min-sup form of Theorem~\ref{thm:minmax} with $W$ spanned by the first $m-1$ eigenvectors of $K$ and the first $n-1$ eigenvectors of $L$.
\end{hint}
\begin{hint}{ex:graphlaplacian}
Expand $\sum_E|x_i-x_j|^2$ and count how often each $|x_i|^2$ appears. For (c), try $x_j=\omega^{jk}$ with $\omega=e^{2\pi i/n}$.
\end{hint}
\begin{hint}{ex:polar}
(a) is one line. For (b), define $U$ on the range of $|T|$ by $U(|T|x)=Tx$, and check that $\ker|T|=\ker T$.
\end{hint}
\begin{hint}{ex:basel}
$G(s,t)=\int_0^1(\mathbf 1_{[0,s]}(u)-s)(\mathbf 1_{[0,t]}(u)-t)\dd u$. The eigenvalues of $\mathcal G$ are the $\frac1{\pi^2n^2}$.
\end{hint}
\begin{hint}{ex:volterranorm}
$V^*V$ has kernel $1-\max(x,y)$, which is $\min(1-x,1-y)$: reflect $[0,1]$.
\end{hint}

\section{Solutions}

\begin{solution}{ex:kl}
\idea A Gaussian process is determined by its covariance, and Mercer's theorem writes the covariance as a sum of products. Each product is the covariance of one random sine wave, and independent amplitudes add
covariances.
\writeup (a) Let $S_N(t)=\sum_{n\le N}\xi_n\sqrt{\lambda_n}e_n(t)$. The $\xi_n$ are orthonormal in $L^2(P)$, so for $M<N$, $E|S_N(t)-S_M(t)|^2=\sum_{M<n\le N}\lambda_ne_n(t)^2$. By Mercer's theorem on the diagonal,
$\sum_n\lambda_ne_n(t)^2=\min(t,t)=t<\infty$, so $(S_N(t))$ is Cauchy in $L^2(P)$, which is complete (Volume~6, Chapter~9).

(b) $e_n(0)=0$ gives $X_0=0$. Inner products are continuous in $L^2(P)$, so by Mercer's theorem
\[E[X_sX_t]=\lim_NE[S_N(s)S_N(t)]=\lim_N\sum_{n\le N}\lambda_ne_n(s)e_n(t)=\min(s,t),\]
and $E[X_t]=0$. $S_N(t)$ is Gaussian
with variance $\sigma_N^2\to t$, and $|E[e^{iuX_t}]-E[e^{iuS_N(t)}]|\le|u|\,E|X_t-S_N(t)|\le|u|\,\|X_t-S_N(t)\|_2\to0$, so $E[e^{iuX_t}]=\lim e^{-u^2\sigma_N^2/2}=e^{-u^2t/2}$: $X_t\sim N(0,t)$. The same argument applied
to any combination $\sum_ia_iX_{t_i}$ shows that $X$ is a Gaussian process with the covariance of Brownian motion.

(c) $G(s,t)=\min(s,t)-st$ is the Green's function of $-u''$ (Lemma~\ref{lem:green} with $q=0$), a positive kernel (Exercise~\ref{ex:basel}), and $\mathcal Ge=\mu e$ with $\mu\neq0$ means $-e''=\mu^{-1}e$ with Dirichlet conditions. By
Theorem~\ref{thm:sturm}, $\mu_n=\frac1{\pi^2n^2}$ and $e_n=\sqrt2\sin(n\pi t)$. The same proof gives the bridge as
\[X_t=\sum_{n\ge1}\xi_n\frac{\sqrt2\sin(n\pi t)}{n\pi},\]
and adding $t\xi_0$ with an independent $\xi_0\sim N(0,1)$ gives Brownian motion again, by the decomposition $B_t=X_t+tZ$ of the Brownian bridge exercise of Volume~3, Chapter~11. The resulting series
$B_t=t\xi_0+\sum_{n\ge1}\xi_n\frac{\sqrt2\sin(n\pi t)}{n\pi}$ is the one Paley and Wiener used in 1934 to construct Brownian motion.
\end{solution}

\begin{solution}{ex:toeplitz}
\idea Convolution commutes with translations, and the characters $e_n$ are the common eigenvectors of all translations. So every convolution operator is diagonal in the Fourier basis, and everything about it is read off
the sequence $\hat g(n)$.
\writeup (a) The kernel $g(x-y)$ is in $L^2([0,1]^2)$ (Tonelli and periodicity), so $T$ is bounded (Example~3.3). With $u=x-y$ and periodicity, $Te_n(x)=\int_0^1g(u)e^{2\pi in(x-u)}\dd u=\hat g(n)e_n(x)$.

(b) As $(e_n)$ is a Hilbert basis of $L^2(\mathbb T)$ (Theorem~2.7 after the change of variable $t=2\pi x$), $Tf=\sum_n\hat g(n)\langle f,e_n\rangle e_n$ and $T^*f=\sum_n\overline{\hat g(n)}\langle f,e_n\rangle e_n$. So $T$ is self-adjoint iff
every $\hat g(n)$ is real, that is iff $g(x)=\overline{g(-x)}$ almost everywhere (the function $x\mapsto\overline{g(-x)}$ has coefficients $\overline{\hat g(n)}$, and Theorem~2.7 identifies a function from its coefficients).
And $\langle Tf,f\rangle=\sum_n\hat g(n)|\langle f,e_n\rangle|^2$ is $\ge0$ for all $f$ iff every $\hat g(n)\ge0$.

(c) The series $\sum r^{|n|}e^{2\pi inx}$ converges absolutely and uniformly, so $P_r$ is continuous, and summing the two geometric series gives $P_r(x)=\frac{1-r^2}{1-2r\cos2\pi x+r^2}$. The kernel $k(x,y)=P_r(x-y)=\sum_nr^{|n|}e_n(x)\overline{e_n(y)}$
is positive in Mercer's sense: $\sum_{i,j}\bar c_ic_jk(x_i,x_j)=\sum_nr^{|n|}\bigl|\sum_jc_j\overline{e_n(x_j)}\bigr|^2\ge0$. Its eigenvalues are the $r^{|n|}$, with sum $\frac{1+r}{1-r}$, and indeed
$\int_0^1k(x,x)\dd x=P_r(0)=\frac{1-r^2}{(1-r)^2}=\frac{1+r}{1-r}$.
\end{solution}

\begin{solution}{ex:weylineq}
\idea The $n$-th eigenvalue is the best upper bound for the Rayleigh quotient once $n-1$ directions are forbidden. Forbid the $m-1$ top directions of $K$ and the $n-1$ top directions of $L$ at once: that is $m+n-2$
directions, and on what remains each operator is small.
\writeup Let $W$ be spanned by the eigenvectors of $K$ for its $m-1$ largest eigenvalues and those of $L$ for its $n-1$ largest (those that exist), so $\dim W\le m+n-2$. For a unit $x\perp W$, the computation in the
proof of Theorem~\ref{thm:minmax} gives $\langle Kx,x\rangle\le\lambda_m(K)$ and $\langle Lx,x\rangle\le\lambda_n(L)$. By the min-sup formula for $K+L$,
\[\lambda_{m+n-1}(K+L)\le\sup_{x\perp W,\,\|x\|=1}\langle(K+L)x,x\rangle\le\lambda_m(K)+\lambda_n(L).\]
If $F\ge0$ has rank $r$, then $\lambda_{r+1}(F)=0$, and the inequality with $L=F$ and $(m,n)$ replaced by $(n,r+1)$ gives $\lambda_{n+r}(K+F)\le\lambda_n(K)$. And $K\le K+F$ gives $\lambda_n(K)\le\lambda_n(K+F)$ by
Corollary~\ref{cor:weylpert}.
\end{solution}

\begin{solution}{ex:graphlaplacian}
\idea The quadratic form of the Laplacian measures how much a function on the vertices varies along the edges. Functions that do not vary at all are constant on components; the cheapest non-constant function
measures how well connected the graph is.
\writeup (a) $\langle Lx,x\rangle=\sum_id_i|x_i|^2-\sum_{i,j}A_{ij}x_j\bar x_i$, and each edge $\{i,j\}$ contributes $|x_i|^2+|x_j|^2-x_j\bar x_i-x_i\bar x_j=|x_i-x_j|^2$. So $L\ge0$, and $Lx=0$ iff $\langle Lx,x\rangle=0$ (for $L\ge0$,
Cauchy--Schwarz for the form $\langle Lx,y\rangle$, as in Proposition~\ref{prop:saspec}, gives $\|Lx\|^2\le\langle Lx,x\rangle^{1/2}\langle L^2x,Lx\rangle^{1/2}$), iff $x_i=x_j$ along every edge, iff $x$ is constant on components.

(b) If $G$ is connected, $\ker L=\C\mathbf 1$, so $\mu_1=0$ is simple. Take an orthonormal eigenbasis $f_1=\frac{\mathbf 1}{\sqrt{|V|}},f_2,\dots$ with eigenvalues $0<\mu_2\le\mu_3\le\cdots$. For $x\perp\mathbf 1$,
$\langle Lx,x\rangle=\sum_{k\ge2}\mu_k|\langle x,f_k\rangle|^2\ge\mu_2\|x\|^2$, with equality at $x=f_2$.

(c) For $C_n$, $(Lx)_j=2x_j-x_{j-1}-x_{j+1}$ (indices mod $n$). With $x_j=\omega^{jk}$, $(Lx)_j=(2-\omega^k-\omega^{-k})x_j=(2-2\cos\frac{2\pi k}n)x_j$. These $n$ vectors are orthogonal (a discrete Fourier basis, Chapter~20), so the
eigenvalues are $2-2\cos\frac{2\pi k}n$, $0\le k<n$, and $\mu_2=2-2\cos\frac{2\pi}n\sim\frac{4\pi^2}{n^2}$. A long cycle is ``almost disconnected'': cut two edges and it falls in two. This is the discrete string of
Theorem~\ref{thm:sturm}, and spectral clustering uses $f_2$ to find such cuts.
\end{solution}

\begin{solution}{ex:polar}
\idea A complex number is $re^{i\theta}$; an operator is ``a stretch, then a rotation''. The stretch $|T|$ is forced by $T^*T$, and the rotation is what is left.
\writeup (a) $\||T|x\|^2=\langle|T|^2x,x\rangle=\langle T^*Tx,x\rangle=\|Tx\|^2$.

(b) By (a), $\ker|T|=\ker T$, and the map $|T|x\mapsto Tx$ is well defined and isometric on $\operatorname{ran}|T|$; it extends by continuity to an isometry $U_0$ of $\overline{\operatorname{ran}|T|}$ onto $\overline{\operatorname{ran}T}$.
Since $|T|$ is self-adjoint, $\overline{\operatorname{ran}|T|}=(\ker|T|)^\perp=(\ker T)^\perp$ (Proposition~3.1). Set $U=U_0$ on $(\ker T)^\perp$ and $U=0$ on $\ker T$. Then $U|T|x=Tx$ for all $x$. Uniqueness: $T=U|T|$ fixes $U$ on
$\operatorname{ran}|T|$, hence on its closure $(\ker T)^\perp$, and $U=0$ on $\ker T$ is imposed.

(c) If $T$ is invertible, $\ker T=\{0\}$ and $\operatorname{ran}T=H$, so $U$ is an isometry of $H$ onto $H$: unitary. For the shift, $S^*S=I$, so $|S|=I$ and $U=S$, an isometry that is not onto. So one cannot always choose
$U$ unitary, unlike in finite dimension where an isometry is automatically onto.
\end{solution}

\begin{solution}{ex:basel}
\idea The string with fixed ends has frequencies $\pi n$, and its Green's function has an elementary diagonal. Mercer's trace formula turns ``sum of the eigenvalues'' into ``integral of the diagonal'': Euler's
series appears as the total compliance of a string.
\writeup (a) For $s,t\in[0,1]$,
\[\int_0^1(\mathbf 1_{[0,s]}-s)(\mathbf 1_{[0,t]}-t)=\min(s,t)-ts-st+st=G(s,t),\]
so $\sum\bar c_ic_jG(t_i,t_j)=\int_0^1\bigl|\sum_jc_j(\mathbf 1_{[0,t_j]}-t_j)\bigr|^2\ge0$, and $G$ is real and symmetric.

(b) By Lemma~\ref{lem:green} with $q=0$, $\mathcal G$ is the inverse of $-u''$ with Dirichlet conditions, and by the proof of Theorem~\ref{thm:sturm} its eigenvalues are $\mu_n=\frac1{\pi^2n^2}$ with eigenfunctions $\sqrt2\sin(n\pi x)$.
Mercer's trace formula gives
\[\sum_{n\ge1}\frac1{\pi^2n^2}=\int_0^1G(x,x)\dd x=\int_0^1(x-x^2)\dd x=\frac16.\]

(c) By Proposition~\ref{prop:hs} in the eigenbasis, $\|\mathcal G\|_{HS}^2=\sum_n\mu_n^2=\sum_n\frac1{\pi^4n^4}$, and it equals $\|G\|_2^2=2\int_0^1\int_0^x y^2(1-x)^2\dd y\dd x=\frac23\int_0^1x^3(1-x)^2\dd x=\frac23\cdot\frac1{60}=\frac1{90}$.
\end{solution}

\begin{solution}{ex:volterranorm}
\idea For a non-normal operator the norm is not read off the eigenvalues, but it is read off the eigenvalues of $V^*V$, which is self-adjoint: $\|V\|^2=\|V^*V\|$. And $V^*V$ turns out to be the covariance of Brownian
motion run backwards.
\writeup By Fubini,
\[\langle Vf,g\rangle=\int_0^1\!\!\int_0^xf(y)\dd y\,\overline{g(x)}\dd x=\int_0^1f(y)\overline{\int_y^1g(x)\dd x}\dd y,\]
so $V^*g(y)=\int_y^1g$. Then
\[V^*Vf(x)=\int_x^1\int_0^tf(s)\dd s\dd t=\int_0^1\bigl(1-\max(x,s)\bigr)f(s)\dd s.\]
The kernel $1-\max(x,s)=\min(1-x,1-s)$ is the kernel of Example~\ref{ex:minkernel} transported by the unitary $Rf(x)=f(1-x)$: $V^*V=RTR^{-1}$. So $V^*V$ has the eigenvalues $\frac1{(n-\frac12)^2\pi^2}$, the largest being
$\frac4{\pi^2}$. By Propositions~3.1 and \ref{prop:saspec}(c), $\|V\|^2=\|V^*V\|=\frac4{\pi^2}$, so $\|V\|=\frac2\pi\approx0.64$, while $r(V)=0$.
\end{solution}

\section*{Sources and further reading}

\begin{itemize}
\item D.~Hilbert, \emph{Grundz\"uge einer allgemeinen Theorie der linearen Integralgleichungen}, Teubner, 1912. The six communications of 1904--1910.
\item L.~A.~Steen, ``Highlights in the history of spectral theory'', \emph{Amer. Math. Monthly} 80 (1973), 359--381. Hilbert, the word ``spectrum'', and quantum mechanics.
\item H.~Weyl, ``Das asymptotische Verteilungsgesetz der Eigenwerte linearer partieller Differentialgleichungen'', \emph{Math. Ann.} 71 (1912), 441--479; W.~Arendt, R.~Nittka, W.~Peter and F.~Steiner, ``Weyl's law'', in
\emph{Mathematical Analysis of Evolution, Information, and Complexity}, Wiley-VCH, 2009. The story of Lorentz's challenge.
\item J.~Mercer, ``Functions of positive and negative type, and their connection with the theory of integral equations'', \emph{Phil. Trans. Roy. Soc. London A} 209 (1909), 415--446.
\item R.~Courant and D.~Hilbert, \emph{Methods of Mathematical Physics}, vol.~1, Interscience, 1953. Min--max, Sturm--Liouville, and the vibrating string.
\item M.~Kac, ``Can one hear the shape of a drum?'', \emph{Amer. Math. Monthly} 73 (1966), 1--23; C.~Gordon, D.~Webb and S.~Wolpert, ``One cannot hear the shape of a drum'', \emph{Bull. Amer. Math. Soc.} 27 (1992), 134--138.
\item G.~Borg, ``Eine Umkehrung der Sturm--Liouvilleschen Eigenwertaufgabe'', \emph{Acta Math.} 78 (1946), 1--96.
\item M.~Reed and B.~Simon, \emph{Methods of Modern Mathematical Physics I: Functional Analysis}, Academic Press, 1980, Chapters~VI--VII; W.~Rudin, \emph{Functional Analysis}, 2nd ed., McGraw--Hill, 1991, Chapter~12.
\item M.~Lo\`eve, \emph{Probability Theory}, Van Nostrand, 1955, \S34 (the Karhunen--Lo\`eve expansion); R.~E.~A.~C.~Paley and N.~Wiener, \emph{Fourier Transforms in the Complex Domain}, AMS, 1934, Chapter~IX.
\end{itemize}
